\documentclass[11pt]{article}
\usepackage[margin=1in]{geometry}
\usepackage{amsmath,amssymb,amsthm,mathtools}
\usepackage{booktabs}
\usepackage{hyperref}

\newtheorem{theorem}{Theorem}
\newtheorem{lemma}[theorem]{Lemma}
\newtheorem{proposition}[theorem]{Proposition}
\newtheorem{corollary}[theorem]{Corollary}
\theoremstyle{definition}
\newtheorem{definition}[theorem]{Definition}
\theoremstyle{remark}
\newtheorem{remark}[theorem]{Remark}

\newcommand{\Z}{\mathbb Z}
\newcommand{\Q}{\mathbb Q}
\newcommand{\conv}{\operatorname{conv}}
\newcommand{\lcm}{\operatorname{lcm}}

\title{Disproof of Conjecture 00000007130:\\
the period of an Ehrhart quasi-polynomial need not be\\ the lcm of the denominators}
\author{jilint777}
\date{2026-10-04}

\begin{document}
\sloppy
\maketitle

\begin{abstract}
Conjecture 00000007130 asserts that the period of the Ehrhart quasi-polynomial of a
rational polytope is the least common multiple of its denominators. This is false.
The triangle $T=\conv\{(0,0),(1,\tfrac12),(2,0)\}$ has vertex denominators with
lcm $2$, yet $\#(tT\cap\Z^2)=\frac{(t+1)(t+2)}{2}$ for every integer $t\ge0$. So its
Ehrhart quasi-polynomial is a polynomial and its period is $1$. This is the smallest case of the
\emph{period collapse} of McAllister and Woods. If ``denominators'' is read as the denominators of
the facet inequalities instead, the triangle $T_2=\conv\{(0,0),(1,0),(0,\tfrac12)\}$ refutes the claim:
its facets are $x\ge0$, $y\ge0$, $x+2y\le1$ (lcm $1$), but its Ehrhart quasi-polynomial has
period exactly $2$. The second counterexample also defeats the weaker reading ``the lcm is
\emph{a} period''. Both results are proved in Lean~4 (core library only), with the polytopes,
lattice-point counts, quasi-polynomials and periods defined from scratch.
\end{abstract}

\section{The conjecture and its readings}

The statement (English version) reads:
\begin{quote}
\emph{Definition:} Integer programming with lattice polytopes. \emph{Conjecture:} The period of the
Ehrhart quasi-polynomial is the least common multiple of denominators, and the explicit form is
the inverse normalization.
\end{quote}
The Chinese version says the same. The conjecture is a conjunction. We refute its first clause:
\begin{equation}\label{eq:clause}
\text{the period of the Ehrhart quasi-polynomial of } P \;=\; \lcm(\text{denominators of } P).
\end{equation}
We do not try to interpret the second clause (``the explicit form is the inverse
normalization''). A conjunction with a false conjunct is false.

\paragraph{Lattice versus rational polytopes.} For a lattice polytope (integer vertices) the Ehrhart
function is a polynomial (Ehrhart's theorem) and every denominator is $1$. Then
\eqref{eq:clause} says ``$1=1$'' and has no content. The words ``quasi-polynomial'',
``period'' and ``denominators'' only mean something for \emph{rational} polytopes, whose vertices
have rational coordinates. So the only reading in which the conjecture says anything is the
standard rational setting: $P\subset\mathbb R^d$ a rational polytope and $L_P(t)=\#(tP\cap\Z^d)$.
We adopt this reading. With vertex or facet denominators, the lattice-only reading is trivially
true, and we state this limitation openly. With coefficient denominators (reading (C) below) the
lattice-only reading fails as well. The standard simplex $\conv\{(0,0),(1,0),(0,1)\}$ has
$L(t)=\frac{(t+1)(t+2)}{2}=\frac12t^2+\frac32t+1$, so its period is $1$ but its coefficient lcm is $2$.

\paragraph{``The period''.} A quasi-polynomial has many periods (every multiple of a period is a
period). With the definite article, ``the period'' is the \emph{minimal} period. This is how it is
used in the literature, e.g.\ McAllister--Woods, \emph{The minimum period of the Ehrhart
quasi-polynomial of a rational polytope} \cite{MW}. If the clause only said ``the lcm is
\emph{a} period'', then in the vertex reading it would be Ehrhart's theorem, which is true
(Remark~\ref{rem:aperiod}). Under the facet reading even this weak form fails
(Section~\ref{sec:facet}).

\paragraph{``Denominators''.} We treat three readings.
\begin{enumerate}
\item[(V)] \emph{Vertex denominators} (standard): the lcm of the denominators of all vertex
coordinates, i.e.\ the least $d\ge1$ such that $dP$ is a lattice polytope. Refuted by $T$
(Section~\ref{sec:T}).
\item[(F)] \emph{Facet denominators}: write each facet inequality as $a\cdot x\le b$ with a
primitive integer normal $a$ and $b\in\Q$ in lowest terms, and take the lcm of the denominators of
the $b$'s. Equivalently, it is the least $d$ such that every facet hyperplane of $dP$ contains a
lattice point. Refuted by $T_2$ (Section~\ref{sec:facet}). This holds for both ``the period'' and
``a period''.
\item[(C)] \emph{Coefficient denominators}: the lcm of the denominators of the coefficients of the
quasi-polynomial. Refuted by $T$ (Remark~\ref{rem:coef}).
\end{enumerate}
Any reading that takes the lcm over vertex \emph{and} facet denominators gives $2$ for $T$, and is
also refuted by $T$.

\paragraph{Integer-programming reading.} In integer programming $T$ is cut out by integer data
$Ax\le b$: $-y\le0$, $-x+2y\le0$, $x+2y\le2$. Its vertex denominators have lcm $2$. Its basis
subdeterminants (the $2\times2$ minors of $A$, i.e.\ the pairs of facet normals) are
$\det\begin{psmallmatrix}0&-1\\-1&2\end{psmallmatrix}=-1$,
$\det\begin{psmallmatrix}0&-1\\1&2\end{psmallmatrix}=1$ and
$\det\begin{psmallmatrix}-1&2\\1&2\end{psmallmatrix}=-4$, so the lcm of their absolute values
is $\lcm(1,1,4)=4$. Both $2$ and $4$ differ from the period $1$.

\section{Definitions}

\begin{definition}
A function $L:\Z_{\ge0}\to\Z$ is a \emph{quasi-polynomial with period $p\ge1$} if there are
polynomials $c_0,\dots,c_{p-1}\in\Q[t]$ with $L(t)=c_{t\bmod p}(t)$ for all $t\ge0$. The
\emph{(minimal) period} of $L$ is the least such $p$. It is unique, and every multiple of a period is a period.
\end{definition}

In Lean the rational polynomials are written over a common denominator: $D\cdot L(t)=P_{t\bmod
p}(t)$ with $D\in\Z_{>0}$ and $P_r\in\Z[t]$. This is equivalent.

For a rational polytope $P$, Ehrhart's theorem \cite{Ehr} says that $L_P(t)=\#(tP\cap\Z^d)$ is a
quasi-polynomial and that the vertex-denominator lcm is a period. McMullen \cite{McM} refined
this coefficient by coefficient. McAllister and Woods \cite{MW} showed that the minimal period can be
strictly smaller (``period collapse''). Their family
$\conv\{(0,0),(1,\frac{D-1}{D}),(D,0)\}$ has denominator $D$ and period $1$, and our $T$ is the
case $D=2$. We give a complete self-contained proof for $T$ below.

\section{\texorpdfstring{The counterexample $T=\conv\{(0,0),(1,\frac12),(2,0)\}$}{The counterexample T}}
\label{sec:T}

The vertex coordinates $0,0,1,\frac12,2,0$ are in lowest terms with denominators
$1,1,1,2,1,1$, so
\[
  \lcm(\text{vertex denominators of } T)=2 .
\]

\begin{lemma}[lattice points of $tT$]\label{lem:HT}
Let $t\ge0$ and $(x,y)\in\Z^2$. Then $(x,y)\in tT$ if and only if
\[
  y\ge0,\qquad 2y\le x,\qquad x+2y\le 2t .
\]
All such points lie in the box $0\le x\le 2t$, $0\le y\le t$.
\end{lemma}

\begin{proof}
$(x,y)\in tT$ means $(x,y)=\nu_0(0,0)+\nu_1(1,\frac12)+\nu_2(2,0)$ with $\nu_i\ge0$ rational and
$\nu_0+\nu_1+\nu_2=t$. (For $t=0$ this gives $\{(0,0)\}=0T$.) The vertices are affinely independent, so
the weights are unique: $y=\nu_1/2$ and $x=\nu_1+2\nu_2$ give
\[
  \nu_1=2y,\qquad \nu_2=\tfrac{x-2y}{2},\qquad \nu_0=t-\nu_1-\nu_2=\tfrac{2t-x-2y}{2}.
\]
Nonnegativity of these is exactly the three inequalities. Conversely, if they hold, these $\nu_i$
work. For the box: $0\le 2y\le x\le 2t-2y\le 2t$ and $4y\le x+2y\le 2t$.
\end{proof}

The three inequalities are the facets $-y\le0$, $-x+2y\le0$, $x+2y\le2$ of $T$, scaled by $t$. Their
normals are primitive and their right-hand sides are integers, so the facet lcm (F) of $T$ is $1$.

\begin{theorem}[period collapse]\label{thm:T}
For every integer $t\ge0$,
\[
  L_T(t)=\#(tT\cap\Z^2)=\frac{(t+1)(t+2)}{2}.
\]
\end{theorem}

\begin{proof}
By Lemma~\ref{lem:HT}, row $y$ ($0\le y\le t$) contains the integers $x$ with $2y\le x\le 2t-2y$.
That is $\max(0,\,2t+1-4y)$ points. If $4y\le 2t$ this is $2t+1-4y$. If $4y>2t$ then $4y\ge 2t+2$, so
the row is empty. Hence $L_T(t)=\sum_{0\le y\le t/2}(2t+1-4y)$.
\begin{itemize}
\item $t=2m$: $\sum_{y=0}^{m}(4m+1-4y)=(m+1)(4m+1)-2m(m+1)=(m+1)(2m+1)=\frac{(t+1)(t+2)}2$.
\item $t=2m+1$: $\sum_{y=0}^{m}(4m+3-4y)=(m+1)(4m+3)-2m(m+1)=(m+1)(2m+3)=\frac{(t+1)(t+2)}2$.
\end{itemize}
(The Lean proof uses the equivalent recurrence $L_T(t+2)=L_T(t)+2t+5$. Row $y+1$ of $(t+2)T$ has as
many points as row $y$ of $tT$, and row $0$ of $(t+2)T$ adds $2t+5$. Together with $L_T(0)=1$ and
$L_T(1)=3$ this gives the formula by induction.)
\end{proof}

\begin{corollary}\label{cor:T}
The Ehrhart quasi-polynomial of $T$ has minimal period $1$, while the lcm of the vertex
denominators of $T$ is $2$. Hence clause \eqref{eq:clause} fails for $T$ in reading (V).
\end{corollary}

\begin{proof}
Theorem~\ref{thm:T} gives period $1$ with $c_0(t)=\frac12t^2+\frac32t+1$, and no period is smaller
than $1$.
\end{proof}

\begin{remark}[``a period'']\label{rem:aperiod}
The number $2$ \emph{is} a period of $L_T$ (take $c_0=c_1$), as Ehrhart's theorem guarantees. So
the counterexample needs ``the period'' to mean the minimal period. That is the meaning of the
definite article and the meaning in \cite{MW}. In Lean this is \texttt{LT\_period\_denLcm}.
\end{remark}

\begin{remark}[coefficient denominators]\label{rem:coef}
The Ehrhart polynomial $\frac12t^2+\frac32t+1$ has coefficient denominators $2,2,1$, with lcm $2\ne1$. So
reading (C) is refuted by $T$ too.
\end{remark}

\begin{remark}
The same happens for every member of the McAllister--Woods family. \texttt{verify.py} checks
$D=2,\dots,6$: denominator $D$, minimal period $1$.
\end{remark}

\section{\texorpdfstring{The facet reading: $T_2=\conv\{(0,0),(1,0),(0,\frac12)\}$}{The facet reading: T2}}
\label{sec:facet}

$T_2$ has facets $-x\le0$, $-y\le0$, $x+2y\le1$. Each has a primitive integer normal and an integer
right-hand side, so the lcm in reading (F) is $1$. Each facet line passes through two vertices. Its vertex
denominators have lcm $2$.

\begin{lemma}\label{lem:HT2}
For $t\ge0$ and $(x,y)\in\Z^2$: $(x,y)\in tT_2$ iff $x\ge0$, $y\ge0$, $x+2y\le t$. Such points
lie in $[0,t]^2$.
\end{lemma}
\begin{proof}
The unique weights are $\nu_1=x$ (for $(1,0)$), $\nu_2=2y$ (for $(0,\frac12)$) and
$\nu_0=t-x-2y$.
\end{proof}

\begin{theorem}\label{thm:T2}
$4L_{T_2}(t)+(t\bmod 2)=(t+2)^2$ for all $t\ge0$. So $L_{T_2}(t)=\frac{(t+2)^2}{4}$ for even
$t$ and $\frac{(t+1)(t+3)}{4}$ for odd $t$, and the minimal period of $L_{T_2}$ is exactly $2$.
\end{theorem}

\begin{proof}
Row $y$ has $\max(0,t+1-2y)$ points, so $L_{T_2}(t)=\sum_{0\le y\le t/2}(t+1-2y)$. For $t=2m$
this is $(m+1)^2$, and for $t=2m+1$ it is $(m+1)(m+2)$. (Lean: $L(t+2)=L(t)+t+3$.) So $2$ is a period.

Suppose $L_{T_2}$ had period $1$, i.e.\ $L_{T_2}(t)=q(t)$ for all $t\ge0$ with $q\in\Q[t]$. Then
$q(t)=(t+2)^2/4$ for infinitely many (even) $t$, so $q=(t+2)^2/4$. But $q(1)=9/4\ne2=L_{T_2}(1)$.
The Lean proof uses an integral version of this argument. Write $D\cdot L_{T_2}=P$ with $D\ge1$ and
$P\in\Z[t]$, and put $a=2D+1$. Then $a\mid P(a)-P(0)=D(L_{T_2}(a)-1)$. Now
$4D(L_{T_2}(a)-1)=D(a^2+4a-1)=Da(a+4)-D$, so $a\mid D$. This contradicts $0<D<a$.
\end{proof}

\begin{corollary}
In reading (F), clause \eqref{eq:clause} fails for $T_2$: the facet lcm is $1$ and the minimal
period is $2$. Since $1$ is not even \emph{a} period, the weak form fails too.
\end{corollary}

\begin{remark}
Each triangle satisfies the clause under the \emph{other} reading: $T$ has facet lcm $1=$ period, and
$T_2$ has vertex lcm $2=$ period. Lean proves both (\texttt{sanity\_T\_facet},
\texttt{sanity\_T2\_vertex}). This shows the formal claims are not vacuous: the predicates can hold. Taken
together, the two triangles refute every reading (V), (F), (C), and the combined vertex-and-facet
lcm.
\end{remark}

\section{Formalization in Lean 4}

The project \texttt{lean4/} (Lean 4.19.0, core library only) defines, from scratch:
\begin{itemize}
\item \texttt{RQ}, \texttt{RPt}, \texttt{Tri}: rational numbers $\mathrm{num}/\mathrm{den}$, rational
points, triangles by vertices. \texttt{Tri.Reduced} means the coordinates are in lowest terms.
\texttt{Tri.denLcm} is the lcm of the six vertex denominators.
\item \texttt{Tri.MemScaled P t x y}: $(x,y)\in tP$. There are $M>0$ and $\mu_i\ge0$ with
$\sum\mu_i=Mt$ and $(x,y)=\sum(\mu_i/M)v_i$, with fractions cleared by $M$ and by the product of the
denominators.
\item \texttt{Facet}: $a_1x+a_2y\le b$ with $b\in\Q$. \texttt{facetDenLcm}, \texttt{Facet.Primitive}
and \texttt{FacetsOf P fs} (each facet is primitive and its line passes through two vertices).
\item \texttt{hcount fs W H t}: the number of $(x,y)\in\Z^2$ with $0\le x\le Wt$, $0\le y\le Ht$ that
satisfy every facet inequality scaled by $t$, computed by explicit finite counting.
\texttt{Describes P fs W H} states that for all $t,x,y$, \texttt{MemScaled P t x y} holds iff
$(x,y)$ is in the box and satisfies the facets. Then \texttt{hcount fs W H} is exactly
$L_P$.
\item \texttt{peval} (integer polynomials as coefficient lists), \texttt{HasPeriod L p} and
\texttt{IsMinPeriod L p}, with \texttt{isMinPeriod\_unique}.
\end{itemize}

\begin{center}
\begin{tabular}{@{}ll@{}}
\toprule
Lean theorem & content\\
\midrule
\texttt{describes\_T}, \texttt{describes\_T2} & Lemmas \ref{lem:HT}, \ref{lem:HT2} (V- and H-description agree)\\
\texttt{T\_denLcm}, \texttt{Tfacets\_denLcm} & vertex lcm $2$, facet lcm $1$ for $T$\\
\texttt{LT\_formula} & $2L_T(t)=(t+1)(t+2)$\\
\texttt{LT\_minPeriod}, \texttt{LT\_not\_minPeriod\_two} & minimal period of $L_T$ is $1$, not $2$\\
\texttt{LT\_period\_denLcm} & $2$ is \emph{a} period of $L_T$\\
\texttt{L2\_formula}, \texttt{L2\_minPeriod} & $4L_{T_2}(t)+(t\bmod2)=(t+2)^2$; minimal period $2$\\
\texttt{conjecture\_00000007130\_false} & $\neg$\,\texttt{VertexClaim}\\
\texttt{conjecture\_00000007130\_false\_facet} & $\neg$\,\texttt{FacetClaim}\\
\texttt{conjecture\_00000007130\_false\_facet\_weak} & $\neg$\,\texttt{FacetClaimWeak} (``a period'')\\
\bottomrule
\end{tabular}
\end{center}

\texttt{VertexClaim} says: for every rational triangle $P$ in lowest terms, and every facet list and
box that describe it, \texttt{IsMinPeriod (hcount fs W H) P.denLcm}. \texttt{FacetClaim} is the same
with \texttt{facetDenLcm fs} and the extra hypothesis \texttt{FacetsOf P fs}. These are claims over
all rational triangles, and triangles are rational polytopes, so their negations refute the
clause for rational polytopes. There is no \texttt{sorry}, no \texttt{native\_decide} and no added axiom.
\texttt{\#print axioms} reports only \texttt{propext}, \texttt{Classical.choice} and
\texttt{Quot.sound}.

\paragraph{Limitations.} Reading (C), the integer-programming subdeterminant reading and the
standard-simplex example are treated only in this report and in \texttt{verify.py}. The Lean
formalization works with triangles and with an explicit facet list plus box. Their
agreement with the vertex description is proved, not assumed (\texttt{describes\_T},
\texttt{describes\_T2}).

\paragraph{Readings that are \emph{not} refuted.} These readings are true, and we do not claim otherwise.
\begin{itemize}
\item The weak vertex form ``the vertex-denominator lcm is \emph{a} period''. This is Ehrhart's
theorem.
\item The lattice-only reading with vertex or facet denominators. There the clause says ``$1=1$''.
\item The reading in which ``denominators'' means the cyclotomic factors (poles at roots of unity) of
the denominator of the reduced Ehrhart series $\sum_t L_P(t)z^t$. The minimal period is the lcm of
the orders of those roots of unity, so this reading is true. For $T$ the series is $1/(1-z)^3$:
only the pole $z=1$, and period $1$.
\end{itemize}

\section{Independent check}

\texttt{verify.py} uses only the Python standard library. It counts lattice points of $tP$ directly
from the \emph{vertices}, using exact barycentric coordinates (Cramer's rule with fractions) on a
generous box. It derives the primitive facet inequalities from the vertices. It determines the
minimal period by exact interpolation of each residue class with polynomials of degree $\le2$, for
all $t\le60$. Results: $L_T(t)=(t+1)(t+2)/2$, with period $1$, vertex lcm $2$, facet lcm $1$, coefficient lcm
$2$. $L_{T_2}$ has period $2$, vertex lcm $2$, facet lcm $1$. The McAllister--Woods triangles with $D=2..6$
all have period $1$. It also computes the basis subdeterminants of $T$ ($1,1,4$, lcm $4$) and checks
that the standard lattice simplex has period $1$ but coefficient lcm $2$.

\begin{thebibliography}{9}
\bibitem{Ehr} E.~Ehrhart, Sur les poly\`edres rationnels homoth\'etiques \`a $n$ dimensions,
\emph{C.~R. Acad. Sci. Paris} 254 (1962), 616--618.
\bibitem{McM} P.~McMullen, Lattice invariant valuations on rational polytopes, \emph{Arch. Math.}
31 (1978/79), 509--516.
\bibitem{MW} T.~B.~McAllister and K.~M.~Woods, The minimum period of the Ehrhart quasi-polynomial of
a rational polytope, \emph{J. Combin. Theory Ser.~A} 109 (2005), 345--352.
\bibitem{BR} M.~Beck and S.~Robins, \emph{Computing the Continuous Discretely}, 2nd ed., Springer, 2015.
\end{thebibliography}

\end{document}
